\chapter{Testing with users}\label{ch:user-testing}
% Draft by Claude, edit freely. Everything in this chapter is a plan.
% The build test is in red on purpose: we still need to discuss it with Florian.

\fillin{This chapter is a plan. None of these tests have been done yet.}

Our goal is a station that anyone can build (Section~\ref{sec:research-question}). We can only
find out if that is true by watching people who are not us try to build one. This chapter
describes that test and the user studies we leave for next semester.

\section{Build test}\label{sec:build-test}

\begin{planned}
\textbf{Idea only. To be discussed with Florian. Most likely next semester, when the design is
stable enough that the guide does not change every week.}

\paragraph{Goal.} Find out if people without experience in electronics can build a working
station from our build guide, and where they get stuck.

\paragraph{Participants.} 3 to 5 people who have not seen the project and have little or no
experience with electronics or 3D printing.

\paragraph{Method.} Each participant gets a kit (printed parts, electronics, tools) and the build
guide. They build the station while thinking aloud, that is, they say what they are doing and
what confuses them. One of us watches and takes notes but only helps when the participant is
stuck. Afterwards we hold a short interview about the hardest parts. We then fix the guide and
test it with the next participant. The full plan is in Appendix~\ref{app:build-test}.

\paragraph{What we measure.}
\begin{itemize}
  \item Time for each step of the guide.
  \item Mistakes (for example a wire on the wrong pin) and whether the participant found them.
  \item How often, and at which step, the participant needed help.
  \item Whether the station works at the end (it takes a picture of the tape).
\end{itemize}

\paragraph{Open questions.} Should participants build the whole station or only parts of it?
Who pays for the kits? Is it safe enough without a supervisor present (everything runs on 5~V,
but the fan and motor can pinch fingers)?
\end{planned}

\section{Why we do not test the web interface}\label{sec:no-ui-test}

The web interface described in Section~\ref{sec:control-software} is a tool for us while we
develop the device. It has many debug functions that only matter for tests and this report. The
software on a finished station will look very different, so a usability test of today's
interface would tell us little. Setting up a finished station is part of the build test instead.

\section{Not done this semester}\label{sec:not-done-users}

Two more studies would answer our questions better, but they need a device that runs for days
and recognises pollen. We did not get that far this semester.

\paragraph{Stations at people's homes.} We would put 3 to 5 stations at people's homes for two to
four weeks. They would write down short notes each week, and we would interview them at the end.
This shows how much upkeep a station needs in practice (tape, placement, Wi-Fi), and whether
people look at the data at all.
\future{Next semester, once the device can run outdoors for days.}

\paragraph{How to show the data.} We would make simple mock-ups of a local pollen map or app and
show them to the people from the interviews (Section~\ref{sec:allergy-interviews}). A cheap
station is less accurate than a reference trap, so it is important that people understand how
certain a number is. We would compare a few ways of showing that.
\future{Next semester, when we have real counts to show.}
